\subsubsection{几何事件分割与计算缓存}
AEQD线段、反投影航迹和原始DEM分别承担平面路径定义、栅格定位和地形查询。对于基地径向距离，使用椭球测地线作为独立核验对象，相关数值算法可参照Karney\cite{karney2013}。显示地图的重采样不进入能耗和高度计算，保证图形美化不会改变模型。

设一条航迹的格边交点参数及两个端点排序为
\begin{equation}
0=u_0<u_1<\cdots<u_M=1.
\label{eq:geometry_breakpoints}
\end{equation}
在相邻事件之间没有新的格边穿越，其所在像元集合不变。因此每个开区间用内部点定位，同时对交点采用闭像元约定，即可覆盖连续路径实际触及的像元。边界和角点接触不能只由两端采样点推断；完整事件分割正是地形核验的重要步骤。

对固定有向航段缓存水平距离、最高地形、巡航海拔、端点作业高度和升降高度。反向航段可以共享距离和最高地形，但升降高度按起终点重新计算。改变箱组而路线不变时复用几何，重新计算载荷能耗；改变访问顺序则读取新的有向航段；中继位置或悬停高度改变后，转场与通信几何均需更新。这种分层复用减少重复地形查询，又保持候选之间相同的物理口径。
